% Section 3.2: Quadratic forms, polarization and signature | Part: 1 | Depends on: sec:vectors, sec:tensors | Status: done | Original lines: 844-881
%!TEX root = ../../../main.tex
\subsection{Quadratic forms, polarization and signature}\label{sec:quadform}
The geometric algebra of Sections~\ref{sec:clifford}--\ref{sec:duality} is built directly on the quadratic form $Q(v)=g(v,v)$
associated with the metric, so we recall the relation between the two and the notation for the signature.
% REWRITTEN: merged with the formal box "metric" above, which is the definition of a metric. Kept here: the
% quadratic form, polarization, the signature (p,q,r) and R^{p,q}. As in the geometric-algebra literature, this
% definition also allows a degenerate g (r>0); the formal box above does not (logged in docs/NOTATION.md).
\begin{definition}[Metric, quadratic form]
Let $g:V\times V\to\K$ be a symmetric bilinear form (a metric as above, but here non-degeneracy is not assumed). It induces the
\emph{quadratic form} $Q(v)=g(v,v)$, and conversely (when $\mathrm{char}\,\K\neq2$) $g$ is
recovered by \emph{polarization}:
\begin{equation}\label{eq:polar}
  g(u,v)=\tfrac12\bigl(Q(u+v)-Q(u)-Q(v)\bigr).
\end{equation}
$g$ is \emph{non-degenerate} if $g(u,\cdot)=0$ implies $u=0$. By Sylvester's law of inertia, over
$\R$ there is an \emph{orthonormal basis} with $g(e_i,e_j)=0$ for $i\neq j$ and
$g(e_i,e_i)\in\{+1,-1,0\}$; the numbers $(p,q,r)$ of $+1,-1,0$ form the \emph{signature}.
We write $\R^{p,q}$ when $r=0$ and $\R^n=\R^{n,0}$ for the Euclidean case.
\end{definition}

% REWRITTEN: the repeated statement of the musical isomorphism (original L861) dropped; see the formal box above.
We shall write $u\cdot v:=g(u,v)$.

\begin{intuition}
The metric is the ruler-and-protractor: it tells you how long vectors are and what angle they make.
Signature $(p,q)$ says that $p$ directions behave like space and $q$ like time (e.g.\ $\R^{1,3}$ for
special relativity). The exterior algebra will need \emph{no} metric; geometric algebra needs one.
\end{intuition}


\begin{pitfall}[three sign conventions for the signature]
(i) Throughout this document the signature $(p,q)$ counts $p$ \emph{positive} and $q$ \emph{negative} eigenvalues,
so that $\R^{p,q}$ and $\Cl(p,q)$ have $p$ basis vectors with $e_i^2=+1$ and $q$ with $e_i^2=-1$; Euclidean space is $(n,0)$.
(ii) General relativity (and Sections~\ref{sec:phys}--\ref{sec:bridge}) uses ``mostly plus'' $\eta=\mathrm{diag}(-1,1,1,1)$, i.e.\ signature
$(3,1)$; (iii) the spacetime algebra of Section~\ref{sec:ga-calculus} uses ``mostly minus'' $\gamma_0^2=+1$, $\gamma_i^2=-1$, i.e.\ $(1,3)$.
Physical statements are the same in the two conventions; only the signs of formulas such as $\vd^2=\pm\Box$ (for the vector
derivative $\vd$ of Section~\ref{sec:ga-calculus}) or
$p^2=\mp m^2$ differ. Always check which one a formula assumes. This box is the reference for signature conventions in the whole document.
\end{pitfall}
